\documentclass[12pt]{article}
\usepackage[utf8]{inputenc}
\usepackage[T1]{fontenc}
\usepackage[margin=1in]{geometry}
\usepackage{setspace}
\usepackage{xcolor}
\usepackage[colorlinks=true, urlcolor=blue!60!black, linkcolor=blue!60!black]{hyperref}
\setlength{\parindent}{0em}
\setlength{\parskip}{0.6em}
\renewcommand{\baselinestretch}{1.0}
\begin{document}
\begin{center}
{\LARGE \textit{Katzenbach v. McClung}, 379 U.S. 294 (1964)\par}
\vspace{0.5em}
{\large Case Brief --- Two-Page Edition\par}
\end{center}

\section{Memory Jogger}
Congress may ban racial discrimination at a local restaurant, Ollie's Barbecue, because a substantial portion of its food moved in interstate commerce and the aggregate effect burdens trade (\href{https://www.courtlistener.com/opinion/106951/katzenbach-v-mcclung/}{CourtListener}).

\section{Detailed Case Facts}
Since 1927, Ollie's Barbecue, a Birmingham family restaurant, had refused dining-room service to African Americans. Of about \$150,000 in annual food purchases, \$69,683 (46\%) was meat from outside Alabama. Title II of the Civil Rights Act of 1964 covered restaurants serving interstate travelers or food a ``substantial portion'' of which had moved in commerce. The government did not claim that interstate travelers ate there.

\section{Procedural History}
A three-judge district court enjoined enforcement, finding no demonstrable link between out-of-state food and discrimination. The Supreme Court reversed on direct appeal, deciding the case with \textit{Heart of Atlanta Motel}.

\section{Judicial Votes}
9--0. Clark wrote for the Court; Black, Douglas, and Goldberg filed concurrences in \textit{Heart of Atlanta} applicable to both cases.

\section{Holding}
Title II validly applies to such a restaurant. Congress may regulate local commercial activity when it has a rational basis for finding that the class, in the aggregate, substantially affects interstate commerce; formal findings and case-by-case proof are unnecessary.

\section{Analysis of Opinions}
\textbf{Majority:} Hearings showed that segregation reduced restaurant purchases of interstate food, depressed business, obstructed travel (``one can hardly travel without eating''), and deterred new industry. Under \textit{Wickard}, Ollie's small purchases, aggregated, were far from trivial; given a rational basis, ``our investigation is at an end.'' \textbf{Black:} whether commerce is affected is ``ultimately a judicial, rather than a legislative, question.'' \textbf{Douglas:} \S~5 of the Fourteenth Amendment would have a ``more settling effect.'' \textbf{Goldberg:} the Act's purpose is ``the vindication of human dignity, and not mere economics.''

\section{Examples: Future Applications}
\textbf{Same-Side:} a discrimination ban reaches local grocery stores buying a quarter of their inventory out of state. \textbf{Opposite-Side:} a ban on possessing knives near parks regulates noneconomic conduct and fails under \textit{Lopez}. \textbf{Fence-Sitter:} a rural farm-to-table diner whose only interstate inputs are salt and a stove.

\section{Critique}
Originalists see \textit{McClung} as the high-water mark of an unbounded Commerce Clause, though \textit{Lopez} preserved it because restaurants are economic. Progressives regret that equal dignity rested on meat shipments rather than the Fourteenth Amendment.

\section{Key Quotations}
\begin{quote}
``Of course, the mere fact that Congress has said when particular activity shall be deemed to affect commerce does not preclude further examination by this Court. But where we find that the legislators, in light of the facts and testimony before them, have a rational basis for finding a chosen regulatory scheme necessary to the protection of commerce, our investigation is at an end.''
\end{quote}
\noindent --- Justice Clark, for the Court, 379 U.S. at 303--304.

\end{document}
